\begin{tikzpicture}[scale=0.92, every node/.style={font=\scriptsize}]
  \draw[fill=softblue,draw=linegray] (0,0) rectangle (6.25,4.45);

  % Panneau gauche : vue dans l'axe du laser.
  \begin{scope}[shift={(0.25,0.15)}]
    \node[heigred,font=\small\bfseries] at (1.25,3.95) {vue laser};
    \draw[fill=white,draw=linegray,rounded corners=2pt] (0.00,0.38) rectangle (2.50,3.70);

    % Plateau.
    \draw[fill=white,draw=mutedgray] (0.28,0.80) -- (1.62,1.05) -- (2.28,0.76) -- (0.95,0.50) -- cycle;

    % Objet ondulé.
    \draw[fill=white,draw=deepblue,thick]
      (1.18,0.88)
      .. controls (1.76,1.17) and (1.10,1.45) .. (1.55,1.78)
      .. controls (2.02,2.12) and (1.22,2.42) .. (1.82,2.78)
      .. controls (2.18,3.03) and (1.82,3.26) .. (2.15,3.42)
      -- (2.28,0.88) -- cycle;
    \draw[deepblue!30,thin] (1.18,0.88) -- (2.28,0.88);
    \draw[deepblue!30,thin] (1.55,1.78) -- (2.28,0.88);
    \draw[deepblue!30,thin] (1.82,2.78) -- (2.28,0.88);

    % Plan laser vu de face : même ligne verticale.
    \draw[fill=softred,draw=heigred,very thick,opacity=0.42] (0.72,0.82) rectangle (1.42,3.40);
    \draw[heigred,very thick] (1.07,0.78) -- (1.07,3.45);
    \node[align=center,text width=2.0cm] at (1.25,0.23) {ligne droite};
  \end{scope}

  % Flèche indiquant le changement de point de vue.
  \draw[->,thick,mutedgray] (2.95,2.55) arc (150:30:0.62);
  \node[mutedgray,font=\small] at (3.12,3.02) {$30^\circ$};

  % Panneau droit : vue caméra décalée.
  \begin{scope}[shift={(3.52,0.15)}]
    \node[deepblue,font=\small\bfseries] at (1.25,3.95) {vue caméra};
    \draw[fill=white,draw=linegray,rounded corners=2pt] (0.00,0.38) rectangle (2.50,3.70);

    % Plateau.
    \draw[fill=white,draw=mutedgray] (0.24,0.80) -- (1.68,1.10) -- (2.30,0.78) -- (0.88,0.48) -- cycle;

    % Plan laser vu de biais.
    \draw[fill=softred,draw=heigred,very thick,opacity=0.34]
      (0.54,0.82) -- (1.28,1.02) -- (1.28,3.42) -- (0.54,3.20) -- cycle;
    \draw[heigred,very thick] (0.54,0.82) -- (0.54,3.20);
    \draw[heigred,very thick] (1.28,1.02) -- (1.28,3.42);

    % Objet ondulé.
    \draw[fill=white,draw=deepblue,thick]
      (1.24,0.88)
      .. controls (1.84,1.18) and (1.20,1.48) .. (1.70,1.82)
      .. controls (2.24,2.18) and (1.44,2.50) .. (2.00,2.84)
      .. controls (2.34,3.08) and (2.00,3.30) .. (2.28,3.42)
      -- (2.38,0.90) -- cycle;
    \draw[deepblue!30,thin] (1.24,0.88) -- (2.38,0.90);
    \draw[deepblue!30,thin] (1.70,1.82) -- (2.38,0.90);
    \draw[deepblue!30,thin] (2.00,2.84) -- (2.38,0.90);

    % Même ligne laser, vue de biais par la caméra.
    \draw[heigred,very thick]
      (0.95,0.82)
      .. controls (1.42,1.18) and (0.82,1.50) .. (1.38,1.92)
      .. controls (1.86,2.30) and (1.06,2.66) .. (1.68,3.05)
      .. controls (1.96,3.23) and (1.62,3.42) .. (1.86,3.58);
    \node[align=center,text width=2.0cm] at (1.25,0.23) {ligne déformée};
  \end{scope}
\end{tikzpicture}
